% Generated draft: requires mathematical and source review.
\documentclass[11pt]{article}
\usepackage[margin=25mm]{geometry}
\usepackage{fontspec}
\setmainfont{Latin Modern Roman}
\usepackage{amsmath,amssymb,mathtools,unicode-math}
\setmathfont{Latin Modern Math}
\usepackage{enumitem,needspace,xurl,hyperref}
\setlength{\parindent}{0pt}
\setlength{\parskip}{5pt}
\setlength{\emergencystretch}{3em}
\newcommand{\N}{\mathbb{N}}
\newcommand{\Z}{\mathbb{Z}}
\newcommand{\Q}{\mathbb{Q}}
\newcommand{\R}{\mathbb{R}}
\newcommand{\C}{\mathbb{C}}
\newcommand{\F}{\mathbb{F}}
\newcommand{\A}{\mathbb{A}}
\newcommand{\PP}{\mathbb{P}}
\newcommand{\E}{\mathbb{E}}
\newcommand{\eps}{\varepsilon}
\newcommand{\dd}{\,\mathrm{d}}
\begin{document}
\section*{UnsolvedMath ID: 10400199}
\textbf{Source problem code:} catalogue entry
\subsection{Definitions and mathematical statement}
A closed oriented \(3\)-manifold \(M\) is an integral homology sphere if \(H_i(M;\Z)\cong H_i(S^3;\Z)\) for all \(i\), and a rational homology sphere if the same holds over \(\Q\). Let \(Z_{\rm LMO}(M)\) denote the Le--Murakami--Ohtsuki invariant, valued in its completed diagram algebra with the source's normalization. Equality of these invariants means equality of the full completed elements, not just a fixed finite-degree truncation.

The literal extracted question asks whether an integral/rational homology \(3\)-sphere \(M\) satisfies
\[
 Z_{\rm LMO}(M)=Z_{\rm LMO}(S^3).
\]
It omits a nontriviality requirement and thus literally permits \(M=S^3\). The supplied research summary explicitly identifies the intended issue as existence of a nontrivial integral homology sphere with this equality, equivalently detection of \(S^3\). Under that intended qualification, the integral version is
\[
 \exists M:\quad H_*(M;\Z)\cong H_*(S^3;\Z),\quad
 M\not\cong S^3,\quad Z_{\rm LMO}(M)=Z_{\rm LMO}(S^3).
\]
The rational version replaces the homology condition by its \(\Q\)-coefficient counterpart. The statement's slash does not decide between these versions, and they are not asserted equivalent. The normalization and degree-zero conventions, particularly for rational homology spheres, must be fixed before the equality is used. The trailing heading ``11.3 Characterization of the image of the LMO invariant'' is surrounding source context, not an additional precise request in this record.
\subsection{Short English statement}
Is there a homology three-sphere other than \(S^3\) with exactly the same full LMO invariant as \(S^3\)? The intended nontriviality comes from the supplied research summary; it is missing from the literal question. Integral and rational versions and normalization must remain distinct.
\subsection{Sources}
\begin{itemize}
\item Local UnsolvedMath registry, permanent ID \(10400199\), exact source code \texttt{catalogue entry}, Problem 11.3. No individual proposer or proposed year is supplied. The extract includes the following Section 11.3 heading on image characterization.
\item Ohtsuki, \emph{Problems on invariants of knots and 3-manifolds} (2002), Problem 11.3, PDF page 156: \url{https://msp.org/gtm/2002/04/gtm-2002-04-024s.pdf}. Bold headings distinguish 231 numbered items from references in surrounding exposition.
\item Extraction pdf-font-structure, accessed 2026-07-29. Source evidence NEEDS\_REVIEW presents an open problem and requires a current primary-literature audit. Registry research\_classification OPEN-TRIAGE and a summary labeled ``Classification: OPEN'' coexist; the summary specifies nontrivial integral homology spheres and calls this the \(S^3\)-detection problem. No further primary citation or resolution is supplied.
\item Dataset attribution: Ulam AI / UnsolvedMath Contributors, CC BY 4.0. Source rights NEEDS\_REVIEW: publicly accessible author/publisher copy, redistribution terms requiring review.
\end{itemize}
\end{document}
